% GENERATED FILE. Edit canonical lessons, then run npm run generate:books.
% canonical-source-hash: fnv1a64:4ab332c7768ea5c7
% canonical-lessons: SA-C254-cakravakah, SA-C254-syenah, SA-C254-pratarasah, SA-C254-ratribhojanam, SA-C254-pakvannam

\chapter{Ruddy goose, Hawk, Breakfast, Dinner, Cooked food}
\label{ch:sa-ruddy-goose-hawk-breakfast-dinner-cooked-food}
\hlchaptermodality{254}

\begin{chapteropening}
\textbf{By the end of this chapter:} \emph{I can say a ruddy goose, a hawk, breakfast, dinner and cooked food.}

Complete the last lesson of chapter 254: I can say a ruddy goose, a hawk, breakfast, dinner and cooked food.
\end{chapteropening}

\section[\texorpdfstring{\sk{चक्रवाकः} (cakravākaḥ)}{चक्रवाकः (cakravākaḥ)}]{\sk{चक्रवाकः} (cakravākaḥ) — a ruddy goose}

\label{lesson:SA-C254-cakravakah}



\begin{quote}
Before the new one: say the Sanskrit for a lizard, then the Sanskrit for a crocodile.
\end{quote}

\subsection*{You'll want to know: \sk{चक्रवाकः}}
\textbf{\sk{चक्रवाकः}} — \emph{cakravākaḥ} — \textquotedblleft{}a ruddy goose\textquotedblright{}.

\begin{grammarlens}[title={What you've built}]
One more word for this chapter.
\end{grammarlens}

\subsection*{Your turn}
\begin{itemize}
\raggedright
  \item \emph{Say it:} \emph{cakravākaḥ}
  \item \emph{Say it:} \emph{cakravākaḥ}, once more
  \item \emph{Say it:} say \emph{cakravākaḥ}
  \item \emph{Recall:} say the Sanskrit for a lizard, then the Sanskrit for a crocodile, then say \emph{cakravākaḥ} again
\end{itemize}

\begin{tcolorbox}[breakable,colback=teal!4,colframe=teal!35!black,title={Before you move on}]
What does \sk{चक्रवाकः} mean? (\textquotedblleft{}A ruddy goose\textquotedblright{}.) Say it once more.
\end{tcolorbox}
\section[\texorpdfstring{\sk{श्येनः} (śyenaḥ)}{श्येनः (śyenaḥ)}]{\sk{श्येनः} (śyenaḥ) — a hawk}

\label{lesson:SA-C254-syenah}



\begin{quote}
Before the new one: say the Sanskrit for a crocodile, then the Sanskrit for a ruddy goose.
\end{quote}

\subsection*{You'll want to know: \sk{श्येनः}}
\textbf{\sk{श्येनः}} — \emph{śyenaḥ} — \textquotedblleft{}a hawk\textquotedblright{}.

\begin{grammarlens}[title={What you've built}]
One more word for this chapter.
\end{grammarlens}

\subsection*{Your turn}
\begin{itemize}
\raggedright
  \item \emph{Say it:} \emph{śyenaḥ}
  \item \emph{Say it:} \emph{śyenaḥ}, once more
  \item \emph{Say it:} say \emph{śyenaḥ}
  \item \emph{Recall:} say the Sanskrit for a crocodile, then the Sanskrit for a ruddy goose, then say \emph{śyenaḥ} again
\end{itemize}

\begin{tcolorbox}[breakable,colback=teal!4,colframe=teal!35!black,title={Before you move on}]
What does \sk{श्येनः} mean? (\textquotedblleft{}A hawk\textquotedblright{}.) Say it once more.
\end{tcolorbox}
\section[\texorpdfstring{\sk{प्रातराशः} (prātarāśaḥ)}{प्रातराशः (prātarāśaḥ)}]{\sk{प्रातराशः} (prātarāśaḥ) — breakfast}

\label{lesson:SA-C254-pratarasah}



\begin{quote}
Before the new one: say the Sanskrit for a ruddy goose, then the Sanskrit for a hawk.
\end{quote}

\subsection*{You'll want to know: \sk{प्रातराशः}}
\textbf{\sk{प्रातराशः}} — \emph{prātarāśaḥ} — \textquotedblleft{}breakfast\textquotedblright{}.

\begin{grammarlens}[title={What you've built}]
One more word for this chapter.
\end{grammarlens}

\subsection*{Your turn}
\begin{itemize}
\raggedright
  \item \emph{Say it:} \emph{prātarāśaḥ}
  \item \emph{Say it:} \emph{prātarāśaḥ}, once more
  \item \emph{Say it:} say \emph{prātarāśaḥ}
  \item \emph{Recall:} say the Sanskrit for a ruddy goose, then the Sanskrit for a hawk, then say \emph{prātarāśaḥ} again
\end{itemize}

\begin{tcolorbox}[breakable,colback=teal!4,colframe=teal!35!black,title={Before you move on}]
What does \sk{प्रातराशः} mean? (\textquotedblleft{}Breakfast\textquotedblright{}.) Say it once more.
\end{tcolorbox}
\section[\texorpdfstring{\sk{रात्रिभोजनम्} (rātribhojanam)}{रात्रिभोजनम् (rātribhojanam)}]{\sk{रात्रिभोजनम्} (rātribhojanam) — dinner}

\label{lesson:SA-C254-ratribhojanam}



\begin{quote}
Before the new one: say the Sanskrit for a hawk, then the Sanskrit for breakfast.
\end{quote}

\subsection*{You'll want to know: \sk{रात्रिभोजनम्}}
\textbf{\sk{रात्रिभोजनम्}} — \emph{rātribhojanam} — \textquotedblleft{}dinner\textquotedblright{}.

\begin{grammarlens}[title={What you've built}]
One more word for this chapter.
\end{grammarlens}

\subsection*{Your turn}
\begin{itemize}
\raggedright
  \item \emph{Say it:} \emph{rātribhojanam}
  \item \emph{Say it:} \emph{rātribhojanam}, once more
  \item \emph{Say it:} say \emph{rātribhojanam}
  \item \emph{Recall:} say the Sanskrit for a hawk, then the Sanskrit for breakfast, then say \emph{rātribhojanam} again
\end{itemize}

\begin{tcolorbox}[breakable,colback=teal!4,colframe=teal!35!black,title={Before you move on}]
What does \sk{रात्रिभोजनम्} mean? (\textquotedblleft{}Dinner\textquotedblright{}.) Say it once more.
\end{tcolorbox}
\section[\texorpdfstring{\sk{पक्वान्नम्} (pakvānnam)}{पक्वान्नम् (pakvānnam)}]{\sk{पक्वान्नम्} (pakvānnam) — cooked food}

\label{lesson:SA-C254-pakvannam}



\begin{quote}
Before the new one: say the Sanskrit for breakfast, then the Sanskrit for dinner.
\end{quote}

\subsection*{You'll want to know: \sk{पक्वान्नम्}}
\textbf{\sk{पक्वान्नम्}} — \emph{pakvānnam} — \textquotedblleft{}cooked food\textquotedblright{}.

\begin{grammarlens}[title={What you've built}]
One more word for this chapter.
\end{grammarlens}

\subsection*{Your turn}
\begin{itemize}
\raggedright
  \item \emph{Say it:} \emph{pakvānnam}
  \item \emph{Say it:} \emph{pakvānnam}, once more
  \item \emph{Say it:} say \emph{pakvānnam}
  \item \emph{Recall:} say the Sanskrit for breakfast, then the Sanskrit for dinner, then say \emph{pakvānnam} again
\end{itemize}

\begin{tcolorbox}[breakable,colback=teal!4,colframe=teal!35!black,title={Before you move on}]
What does \sk{पक्वान्नम्} mean? (\textquotedblleft{}Cooked food\textquotedblright{}.) Say it once more.
\end{tcolorbox}
